\documentclass[10pt,a4paper]{amsart}
\usepackage{amsmath,amssymb,amsthm,graphicx}
\usepackage{verbatim}
\newenvironment{reference}{\comment}{\endcomment}
%\newenvironment{reference}{}{}
\newenvironment{slogan}{\comment}{\endcomment}
\newenvironment{history}{\comment}{\endcomment}
% For commutative diagrams we use Xy-pic
\usepackage[all]{xy}

% We use 2cell for 2-commutative diagrams.
\xyoption{2cell}
% Theorem environments.
%
\theoremstyle{plain}
\newtheorem{theorem}[subsection]{Theorem}
\newtheorem{proposition}[subsection]{Proposition}
\newtheorem{lemma}[subsection]{Lemma}

\theoremstyle{definition}
\newtheorem{definition}[subsection]{Definition}
\newtheorem{example}[subsection]{Example}
\newtheorem{exercise}[subsection]{Exercise}
\newtheorem{situation}[subsection]{Situation}

\theoremstyle{remark}
\newtheorem{remark}[subsection]{Remark}
\newtheorem{remarks}[subsection]{Remarks}

\numberwithin{equation}{subsection}

% Macros
%
\def\lim{\mathop{\mathrm{lim}}\nolimits}
\def\colim{\mathop{\mathrm{colim}}\nolimits}
\def\Spec{\mathop{\mathrm{Spec}}}
\def\Hom{\mathop{\mathrm{Hom}}\nolimits}
\def\Ext{\mathop{\mathrm{Ext}}\nolimits}
\def\SheafHom{\mathop{\mathcal{H}\!\mathit{om}}\nolimits}
\def\SheafExt{\mathop{\mathcal{E}\!\mathit{xt}}\nolimits}
\def\Sch{\mathit{Sch}}
\def\Mor{\mathop{\mathrm{Mor}}\nolimits}
\def\Ob{\mathop{\mathrm{Ob}}\nolimits}
\def\Sh{\mathop{\mathit{Sh}}\nolimits}
\def\NL{\mathop{N\!L}\nolimits}
\def\CH{\mathop{\mathrm{CH}}\nolimits}
\def\proetale{{pro\text{-}\acute{e}tale}}
\def\etale{{\acute{e}tale}}
\def\QCoh{\mathit{QCoh}}
\def\Ker{\mathop{\mathrm{Ker}}}
\def\Im{\mathop{\mathrm{Im}}}
\def\Coker{\mathop{\mathrm{Coker}}}
\def\Coim{\mathop{\mathrm{Coim}}}

% Boxtimes
%
\DeclareMathSymbol{\boxtimes}{\mathbin}{AMSa}{"02}

%
% Macros for moduli stacks/spaces
%
\def\QCohstack{\mathcal{QC}\!\mathit{oh}}
\def\Cohstack{\mathcal{C}\!\mathit{oh}}
\def\Spacesstack{\mathcal{S}\!\mathit{paces}}
\def\Quotfunctor{\mathrm{Quot}}
\def\Hilbfunctor{\mathrm{Hilb}}
\def\Curvesstack{\mathcal{C}\!\mathit{urves}}
\def\Polarizedstack{\mathcal{P}\!\mathit{olarized}}
\def\Complexesstack{\mathcal{C}\!\mathit{omplexes}}
% \Pic is the operator that assigns to X its picard group, usage \Pic(X)
% \Picardstack_{X/B} denotes the Picard stack of X over B
% \Picardfunctor_{X/B} denotes the Picard functor of X over B
\def\Pic{\mathop{\mathrm{Pic}}\nolimits}
\def\Picardstack{\mathcal{P}\!\mathit{ic}}
\def\Picardfunctor{\mathrm{Pic}}
\def\Deformationcategory{\mathcal{D}\!\mathit{ef}}
\usepackage{cleveref}
\crefname{originalreference}{original reference}{original references}
\Crefname{originalreference}{Original reference}{Original references}
\makeatletter
\newlabel{coherent-lemma-coherent-on-proj}{{coherent-lemma-coherent-on-proj}{source}{}{sourcepointer-1}{}}
\newlabel{coherent-lemma-coherent-on-proj@cref}{{[originalreference][0][]coherent-lemma-coherent-on-proj}{[1][0][]source}{}{}{}}
\makeatother
\crefname{originalreference}{original reference}{original references}
\Crefname{originalreference}{Original reference}{Original references}
\makeatletter
\newlabel{coherent-lemma-recover-tail-graded-module}{{coherent-lemma-recover-tail-graded-module}{source}{}{sourcepointer-2}{}}
\newlabel{coherent-lemma-recover-tail-graded-module@cref}{{[originalreference][0][]coherent-lemma-recover-tail-graded-module}{[1][0][]source}{}{}{}}
\makeatother
\crefname{originalreference}{original reference}{original references}
\Crefname{originalreference}{Original reference}{Original references}
\makeatletter
\newlabel{perfect-lemma-acyclicity-lemma-global}{{perfect-lemma-acyclicity-lemma-global}{source}{}{sourcepointer-3}{}}
\newlabel{perfect-lemma-acyclicity-lemma-global@cref}{{[originalreference][0][]perfect-lemma-acyclicity-lemma-global}{[1][0][]source}{}{}{}}
\makeatother
\title{The derived pushforward of the twist of a pulled back quasi-coherent module on the blowup computes the derived tensor product with the ideal power}
\author{The Stacks Project Authors}
\date{stacks-project:local-cohomology.tex:local-snapshot-20261012; GNU FDL 1.2 or later}
\begin{document}
\maketitle
\noindent\textit{Editorial note.} All delivered mathematical text is the author's own, copied from the pinned Stacks Project chapter source sources/stacks-local-cohomology/pinned-original-source.tex; no statement or proof step was rewritten, reordered or removed.\par\medskip
\noindent\textit{Editorial note.} Preamble provenance: the Stacks Project keeps one single shared preamble for every chapter; this card takes it directly from the staged sources/stacks-local-cohomology/preamble.tex as three verbatim spans (lines 8-12, 14-18 and 154-217), so no derived file was written and nothing was added or redefined. The wrapper class and the base packages are supplied by the assembler.\par\medskip
\noindent\textit{Editorial note.} Bibliography provenance: neither the counted statement nor its proof contains a citation (this lane rescans the whole span before writing the card), so the card carries no bibliography and no entry was reconstructed.\par\medskip
\noindent\textit{Editorial note.} Selection provenance: this card was opened under the file-based control contract (CONTROL.md), with the mechanical pre-card gate run on its statement and proof interval before the card was written.\par\medskip
\noindent\textit{Allowed standard prerequisites (editorial).}\begin{itemize}
\item The blowing up p : X \textbackslash{}to Spec(A) of a Noetherian ring in an ideal I as the Proj of the Rees algebra, and the threshold q(A, I) beyond which the twists of the structure sheaf have only the stated cohomology. The author's setup paragraph, reproduced verbatim as a context span of this card (source lines 4604-4616). Fixes all objects and the numerical hypothesis appearing in the counted statement.
\item The derived pullback Lp\textasciicircum{}* of quasi-coherent modules and the fact that a flat resolution of the module computes it. Used inside the counted proof at source lines 4627-4632. Replaces the derived pullback by an explicit complex of pulled back free modules.
\item The acyclicity lemma for the pushforward along the projection, in the form recorded in the chapter on derived categories of schemes. Invoked inside the counted proof at source lines 4636-4637 as a cross-chapter pointer. Shows that the terms of the complex are acyclic for the global sections functor, so the total derived functor is computed termwise.
\end{itemize}\medskip
\ifdefined\hypertarget\hypertarget{sourcepointer-1}{}\fi
\noindent\textit{Source reference.} Original reference coherent-lemma-coherent-on-proj: sources/stacks-local-cohomology/pinned-original-source.tex, source lines 4612--4612. Navigation pointer only: the counted text refers to the Stacks item coherent-lemma-coherent-on-proj, which belongs to another chapter of the project and is not part of this card. The pointer carries the target's own label as its number because this chapter's staged snapshot does not contain the other chapter's tag table; its content is not reproduced here.\par\medskip
\ifdefined\hypertarget\hypertarget{sourcepointer-2}{}\fi
\noindent\textit{Source reference.} Original reference coherent-lemma-recover-tail-graded-module: sources/stacks-local-cohomology/pinned-original-source.tex, source lines 4613--4613. Navigation pointer only: the counted text refers to the Stacks item coherent-lemma-recover-tail-graded-module, which belongs to another chapter of the project and is not part of this card. The pointer carries the target's own label as its number because this chapter's staged snapshot does not contain the other chapter's tag table; its content is not reproduced here.\par\medskip
\ifdefined\hypertarget\hypertarget{sourcepointer-3}{}\fi
\noindent\textit{Source reference.} Original reference perfect-lemma-acyclicity-lemma-global: sources/stacks-local-cohomology/pinned-original-source.tex, source lines 4637--4637. Navigation pointer only: the counted text refers to the Stacks item perfect-lemma-acyclicity-lemma-global, which belongs to another chapter of the project and is not part of this card. The pointer carries the target's own label as its number because this chapter's staged snapshot does not contain the other chapter's tag table; its content is not reproduced here.\par\medskip
\begin{lemma}
\label{lemma-compute-tor-Iq}
In the situation above, for $q \geq q(A, I)$ and any $A$-module $M$ we have
$$
R\Gamma(X, Lp^*\widetilde{M}(q)) \cong M \otimes_A^\mathbf{L} I^q
$$
in $D(A)$.
\end{lemma}

\noindent
For the rest of the discussion in this section we fix a Noetherian ring $A$
and an ideal $I \subset A$. We denote
$$
p : X \to \Spec(A)
$$
the blowing up of $\Spec(A)$ in the ideal $I$. In other words, $X$ is the
$\text{Proj}$ of the Rees algebra $\bigoplus_{n \geq 0} I^n$.
By Cohomology of Schemes, Lemmas \ref{coherent-lemma-coherent-on-proj} and
\ref{coherent-lemma-recover-tail-graded-module}
we can choose an integer $q(A, I) \geq 0$ such that for all $q \geq q(A, I)$
we have $H^i(X, \mathcal{O}_X(q)) = 0$ for
$i > 0$ and $H^0(X, \mathcal{O}_X(q)) = I^q$.

\begin{proof}
Choose a free resolution $F_\bullet \to M$. Then $\widetilde{F}_\bullet$
is a flat resolution of $\widetilde{M}$. Hence $Lp^*\widetilde{M}$
is given by the complex $p^*\widetilde{F}_\bullet$. Thus
$Lp^*\widetilde{M}(q)$ is given by the complex $p^*\widetilde{F}_\bullet(q)$.
Since $p^*\widetilde{F}_i(q)$ are right acyclic for $\Gamma(X, -)$ by
our choice of $q \geq q(A, I)$ and since we have
$\Gamma(X, p^*\widetilde{F}_i(q)) = I^qF_i$ by our choice of $q \geq q(A, I)$,
we get that $R\Gamma(X, Lp^*\widetilde{M}(q))$ is given by the complex
with terms $I^qF_i$ by Derived Categories of Schemes, Lemma
\ref{perfect-lemma-acyclicity-lemma-global}.
The result follows as the complex $I^qF_\bullet$ computes
$M \otimes_A^\mathbf{L} I^q$ by definition.
\end{proof}

\end{document}
